% ══════════════════════════════════════════════════════════════════════════════
%  EXPERIMENT 4 — RATES OF EVAPORATION                           (Lesson 11, Term 2 Week 4)
% ══════════════════════════════════════════════════════════════════════════════
\experimentdiv{Rates of Evaporation}{Writing an aim \textbullet{} Planning your own method \textbullet{} Measuring mass loss over time}{Lesson 11 \textbullet\ Term 2, Week 4 \textbullet\ Lab}\label{sec:4}
\lessoncodes{VC2S8U05 \textbullet{} VC2S8I01--I04, I06}

\vocabhead
\begin{ctab}
\begin{tabular}{|SW{4.2cm}|W{11.6cm}|}
  \hline\hd & \thead{Word} & \thead{Definition (copy from the slide)} \\ \hline
  & \scaf{\textbf{Aim}} & \scaf{A sentence saying what the experiment is trying to find out.} \\ \hline
  & \textbf{Rate of evaporation} & \\ \hline
  & \textbf{Method} & \\ \hline
  & \textbf{Repeatable} & \\ \hline
  & \textbf{Mass lost} & \\ \hline
\end{tabular}
\end{ctab}

\begin{lipanel}
\inquiryli{Write an aim and a hypothesis, then plan and follow a repeatable method.}
\understandli{Explain how one factor affects the rate of evaporation.}
\sctitle
\begin{itemize}[leftmargin=1.4em, itemsep=0pt, topsep=1pt]
  \item I can write an aim and a hypothesis for my own investigation.
  \item I can write numbered method steps that someone else could follow.
  \item I can calculate mass lost and use it to compare rates.
\end{itemize}
\end{lipanel}

\begin{skillsfocus}
\focus{QP2}{Write an aim for your investigation.}
\focus{QP1}{Write a hypothesis with a reason.}
\focus{PC1}{Plan a method someone else could repeat.}
\focus{PC2}{Choose the variable to change and those to keep the same.}
\focus{PC4}{Weigh to 0.1\,g on the same balance each time.}
\focus{PA2}{Graph the mass lost.}
\focus{PA3}{Calculate the mass lost.}
\focus{EV3}{Find sources of error.}
\end{skillsfocus}

% ─── TECHNIQUE ────────────────────────────────────────────────────────────────
\exppart{Technique spotlight}{Blue}

\begin{technique}{Writing an aim}{QP2}
An aim names the independent and dependent variables and how they are linked:
\begin{center}\emph{To find out how} [the independent variable] \emph{affects} [the dependent variable].\end{center}
\end{technique}

\begin{workedex}{An aim}
\emph{To find out how the temperature of water affects the time sugar takes to dissolve.}
\end{workedex}

\begin{wedo}{An aim}
Write an aim for an experiment that tests whether a fan changes how quickly a wet towel dries.
\answerbox{3cm}
\end{wedo}

\begin{technique}{Writing a repeatable method}{PC1}
A method is \textbf{repeatable} when another group could follow it and get similar results. Number each step, start each step with a verb (\emph{Measure}, \emph{Place}, \emph{Record}), and give every amount with its unit.
\end{technique}

\begin{workedex}{Vague step and precise step}
\textbf{Vague:} \emph{Put some water in the dish.}\qquad
\textbf{Precise:} \emph{Measure 20\,mL of water and pour it into the Petri dish.}
\end{workedex}

\begin{youdobox}[Precise steps]
\task Rewrite each step so that it is precise. \skr{PC1}
\pt{a}Weigh it.
\answerbox{2cm}
\pt{b}Leave it somewhere for a while.
\answerbox{2cm}
\end{youdobox}

% ─── PLAN ────────────────────────────────────────────────────────────────────
\exppart{Plan your investigation}{Teal}

Choose \textbf{one} factor to test. Tick it.

\medskip
\begin{tabular}{@{}l l@{}}
\tickbox\enspace Surface area (Petri dish, beaker, test tube) & \tickbox\enspace Air movement (fan on, fan off) \\[3pt]
\tickbox\enspace Temperature (warm spot, cool spot) & \tickbox\enspace Covered or uncovered \\
\end{tabular}

\task Write your aim. \skr{QP2}
\answerbox{3cm}

\needspace{8\baselineskip}
\task Write your hypothesis. \skr{QP1}
\answerbox{3.5cm}

\needspace{10\baselineskip}
\task Name your variables. \skr{PC2}
\begin{ctab}
\begin{tabular}{|SW{4.4cm}|W{10.8cm}|}
  \hline\hd & \thead{Variable} & \thead{In my investigation} \\ \hline
  & \scaf{Dependent} & \scaf{Mass of water lost (g)} \\ \hline
  & Independent & \\ \hline
  & Controlled & \\ \hline
  & Controlled & \\ \hline
  & Controlled & \\ \hline
\end{tabular}
\end{ctab}

\subsubsection*{Equipment}
Containers for your factor \textbullet{} 50\,mL measuring cylinder \textbullet{} electronic balance (0.1\,g) \textbullet{} water \textbullet{} marker and masking tape \textbullet{} anything else your factor needs (fan, plastic wrap, a warm and a cool place)

\needspace{14\baselineskip}
\task Write your method as numbered steps. Include how you will weigh the containers and how often. \skr{PC1}
\answerbox{9cm}

\begin{infobox}{Safety}
Keep the water well away from the balance's power cable and from any fan. Wipe up spills straight away. Have your teacher check your method before you start.\par\smallskip
Teacher's initials: \rule{2.5cm}{0.4pt}
\end{infobox}

% ─── RESULTS ─────────────────────────────────────────────────────────────────
\exppart{Results}{Amber}

\subsubsection*{Table 4.1: Mass of each container and its water \skr{PC4}}
\begin{ctab}
\begin{tabular}{|SW{2.8cm}|C{2.2cm}|C{2.2cm}|C{2.2cm}|C{2.2cm}|C{2.2cm}|}
  \hline\hd & \thead{Container or condition} & \thead{Start (g)} & \thead{Lesson 2 (g)} & \thead{Lesson 3 (g)} & \thead{Lesson 4 (g)} & \thead{Mass lost (g)} \\ \hline
  & \scaf{Petri dish} & \scaf{62.4} & \scaf{58.1} & \scaf{53.7} & \scaf{49.5} & \scaf{12.9} \\ \hline
  & & & & & & \\ \hline
  & & & & & & \\ \hline
  & & & & & & \\ \hline
\end{tabular}
\end{ctab}

\task Show how you calculated one mass lost. \skr{PA3}
\workbox{2.5cm}

\needspace{18\baselineskip}
\task Draw a column graph of the mass lost by each container or condition. Label both axes and give the graph a title. \skr{PA2}
\graphgrid{9cm}

% ─── ANALYSIS ────────────────────────────────────────────────────────────────
\exppart{Analysis and evaluation}{Coral}

\task Which container or condition had the fastest rate of evaporation? How do you know? \skr{PA3}
\answerbox{3.5cm}

\needspace{10\baselineskip}
\task Do your results support your hypothesis? Use numbers from Table 4.1. \skr{EV1}
\answerbox{4cm}

\needspace{12\baselineskip}
\task Explain your results using the particle model.
\answerbox{5cm}

\needspace{12\baselineskip}
\task Give two sources of error and how each could be reduced. \skr{EV3}
\begin{ctab}
\begin{tabular}{|SW{7.2cm}|W{7.8cm}|}
  \hline\hd & \thead{Source of error} & \thead{How to reduce it} \\ \hline
  & \scaf{Water spilled when moving the container} & \scaf{Weigh the container where it stands, or carry it on a tray} \\ \hline
  & & \\ \hline
  & & \\ \hline
\end{tabular}
\end{ctab}

\needspace{10\baselineskip}
\task Was your method repeatable? Name one step another group found hard to follow, and rewrite it. \skr{PC1}
\answerbox{4cm}

\begin{yr8box}
\begin{keepext}
\exttask Rate means \emph{how much per unit of time}. Calculate the rate of evaporation of your fastest container in grams per day, and use it to predict how many days it would take for all its water to evaporate.
\workbox{5cm}
\end{keepext}
\end{yr8box}

% ─── SKILL CHECK ─────────────────────────────────────────────────────────────
\exppart{Skill check}{NavyBlue}

\begin{skillcheck}
\scrow{QP2}{I can write an aim that names both variables.}
\scrow{QP1}{I can write a hypothesis with a reason.}
\scrow{PC1}{I wrote a method another group could repeat.}
\scrow{PC2}{I changed one variable and kept the others the same.}
\scrow{PC4}{I weighed to 0.1\,g on the same balance each time.}
\scrow{PA2}{I drew a column graph of my results.}
\scrow{PA3}{I calculated mass lost correctly.}
\scrow{EV3}{I can name sources of error and how to reduce them.}
\end{skillcheck}

\task Which skill will you focus on in Experiment 5? Explain what you will do differently.
\answerbox{5cm}
